\chapter{المتغيرات العشوائية المتصلة}\label{ch:g12:contdist}

أزمنة الانتظار، والقياسات الفيزيائية، والنسب: كثير من المقادير العشوائية
تأخذ متصلًا من القيم، ولا قيمة مفردة لها \omterm{def:g10:proba:distribution}{احتمال} موجب.
وعندئذٍ تُحسب الاحتمالات \omterm{def:g12:integ:area}{بتكامل} \emph{\omterm{def:g12:contdist:density}{كثافة}}. ويدرس هذا
الفصل التوزيعات المنتظم والأسي والطبيعي، ويستعمل
القانون الطبيعي لتكميم تذبذبات الاستطلاعات والعينات.

\section{كثافات الاحتمال}

\begin{definition}[الكثافة والمتغير العشوائي المتصل]\label{def:g12:contdist:density}
\emph{كثافة الاحتمال}\index{كثافة} على \omterm{def:g10:numbers:interval}{فترة} $I$ هي
\omterm{def:g12:limcont:continuity}{دالة متصلة} غير سالبة $f$ على $I$ حيث $\int_I f(t)\,\dd t = 1$
(ويُفهم \omterm{def:g12:integ:area}{التكامل} على $I$ غير \omterm{def:g12:seq:bounded}{المحدودة} نهايةً لتكاملات
على فترات \omterm{def:g12:seq:bounded}{محدودة} متنامية). ويكون \omterm{def:g12:randvar:rv}{للمتغير العشوائي} $X$
الكثافة $f$ إذا كان من أجل كل $a \leq b$ في $I$:
\[
\P(a \leq X \leq b) = \int_a^b f(t)\,\dd t .
\]
\end{definition}

فالاحتمالات \emph{مساحات تحت \omterm{def:g10:functions:graph}{منحنى} \omterm{def:g12:contdist:density}{الكثافة}}. وخاصة
$\P(X = a) = \int_a^a f = 0$ من أجل كل قيمة مفردة $a$: فلا تحمل \omterm{def:g10:proba:distribution}{الاحتمال} إلا
الفترات، ويكون $\P(a \leq X \leq b) = \P(a < X < b)$.

\begin{definition}[الأمل الرياضي والتباين]\label{def:g12:contdist:exp}
من أجل $X$ ذي \omterm{def:g12:contdist:density}{الكثافة} $f$ على $I$:
\[
\E(X) = \int_I t\,f(t)\,\dd t, \qquad
\V(X) = \int_I \bigl(t - \E(X)\bigr)^2 f(t)\,\dd t
= \E(X^2) - \E(X)^2 .
\]
\end{definition}

وتبقى قواعد الفصلين 14--15 (خطية \omterm{def:g12:randvar:exp}{الأمل الرياضي}، وجمعية
\omterm{def:g12:randvar:exp}{التباين} من أجل المتغيرات المستقلة، ومتراجحة بيينيمي--تشيبيشيف) صالحة؛
وبراهينها، بالتكاملات محل المجاميع، مقبولة في هذا المستوى.

\section{التوزيع المنتظم}

\begin{definition}[التوزيع المنتظم]\label{def:g12:contdist:uniform}
يتبع $X$ \emph{التوزيع المنتظم}\index{توزيع منتظم}
$\mathcal U\intcc ab$ إذا كانت كثافته ثابتة،
$f(t) = \frac{1}{b-a}$ على $\intcc{a}{b}$. وعندئذٍ يكون من أجل
$\intcc{c}{d} \subseteq \intcc ab$،
$\P(c \leq X \leq d) = \frac{d - c}{b - a}$: أي \omterm{def:g10:proba:distribution}{احتمال} يتناسب مع
الطول.
\end{definition}

\begin{proposition}\label{prop:g12:contdist:uniformmoments}
إذا كان $X \sim \mathcal U\intcc ab$، فإن
$\E(X) = \dfrac{a+b}{2}$ و $\V(X) = \dfrac{(b-a)^2}{12}$.
\end{proposition}

\begin{proof}
$\E(X) = \int_a^b \frac{t}{b-a}\dd t
= \frac{b^2 - a^2}{2(b-a)} = \frac{a+b}{2}$. وأما \omterm{def:g12:randvar:exp}{التباين}، فإن
$\E(X^2) = \frac{b^3 - a^3}{3(b-a)} = \frac{a^2 + ab + b^2}{3}$، و
\[
\V(X) = \frac{a^2 + ab + b^2}{3} - \frac{(a+b)^2}{4}
= \frac{4a^2 + 4ab + 4b^2 - 3a^2 - 6ab - 3b^2}{12}
= \frac{(b - a)^2}{12}. \qedhere
\]
\end{proof}

\section{التوزيع الأسي}

\begin{definition}[التوزيع الأسي]\label{def:g12:contdist:expo}
من أجل $\lambda > 0$، يتبع $X$ \emph{التوزيع
الأسي}\index{توزيع أسي} $\mathcal E(\lambda)$ إذا كانت
كثافته على $\intco{0}{+\infty}$ هي
\[
f(t) = \lambda\,\eu^{-\lambda t}.
\]
وعندئذٍ $\P(X \leq x) = 1 - \eu^{-\lambda x}$ و
$\P(X > x) = \eu^{-\lambda x}$ من أجل $x \geq 0$.
\end{definition}

\begin{proposition}\label{prop:g12:contdist:expomoments}
إذا كان $X \sim \mathcal E(\lambda)$:
$\E(X) = \dfrac{1}{\lambda}$ و $\V(X) = \dfrac{1}{\lambda^2}$.
\end{proposition}

\begin{proof}
\omterm{thm:g12:integ:ibp}{بالتكامل بالتجزئة} على $\intcc{0}{A}$:
\[
\int_0^A t\,\lambda\eu^{-\lambda t}\dd t
= \bigl[-t\,\eu^{-\lambda t}\bigr]_0^A + \int_0^A \eu^{-\lambda t}\dd t
= -A\eu^{-\lambda A} + \frac{1 - \eu^{-\lambda A}}{\lambda}
\xrightarrow[A\to+\infty]{} \frac1\lambda,
\]
باستعمال $A\eu^{-\lambda A} \to 0$ (\cref{thm:g12:exp:growth}). ويعطي \omterm{def:g12:integ:area}{تكامل}
ثانٍ بالتجزئة أن $\E(X^2) = \frac{2}{\lambda^2}$، ومنه
$\V(X) = \frac{2}{\lambda^2} - \frac{1}{\lambda^2}$.
\end{proof}

\begin{theorem}[انعدام الذاكرة]\label{thm:g12:contdist:memoryless}
\index{انعدام الذاكرة}
إذا كان $X \sim \mathcal E(\lambda)$، فإنه من أجل كل $s, t \geq 0$:
\[
\pcond{X > s}{X > s + t} = \P(X > t) .
\]
\omterm{def:g12:contdist:expo}{فالتوزيع الأسي} هو قانون الأعمار \emph{بلا شيخوخة}
(النوى المشعّة، لا المصابيح الكهربائية).
\end{theorem}

\begin{proof}
\[
\pcond{X > s}{X > s+t}
= \frac{\P(X > s + t)}{\P(X > s)}
= \frac{\eu^{-\lambda(s+t)}}{\eu^{-\lambda s}} = \eu^{-\lambda t}
= \P(X > t). \qedhere
\]
\end{proof}

\begin{omfigure}
\begin{tikzpicture}
\begin{axis}[omaxis,
  width=8.8cm, height=5cm,
  xmin=-0.4, xmax=4.4, ymin=0, ymax=1.15,
  xtick={1.3}, xticklabels={$t$}, ytick={1}, yticklabels={$\lambda$}]
  \addplot[name path=dens, omDef, thick, domain=0:4.2, samples=80]
    {exp(-x)};
  \path[name path=xaxis] (0,0) -- (4.2,0);
  \addplot[omThm!20] fill between[of=dens and xaxis,
    soft clip={domain=1.3:4.2}];
  \draw[black, dashed, thin] (axis cs:1.3,0) -- (axis cs:1.3,0.2725);
  \node[omThm, font=\small, anchor=west] at (axis cs:2.1,0.35)
    {$\P(X > t) = \eu^{-\lambda t}$};
\end{axis}
\end{tikzpicture}

{\small \omterm{def:g12:contdist:density}{الكثافة} الأسية $\lambda\eu^{-\lambda x}$ (وهنا
$\lambda = 1$): مساحة الذيل وراء $t$ (بالأحمر) هي $\eu^{-\lambda t}$،
ويقول انعدام الذاكرة إن كل ذيل يبدو \omterm{def:g12:randvar:rv}{كالتوزيع} كله
مكبَّرًا.}
\end{omfigure}

\section{التوزيع الطبيعي}

\begin{definition}[التوزيع الطبيعي]\label{def:g12:contdist:normal}
يتبع $X$ \emph{التوزيع الطبيعي المعياري}\index{توزيع
طبيعي} $\mathcal N(0, 1)$ إذا كانت كثافته على $\R$ هي
\[
\varphi(t) = \frac{1}{\sqrt{2\pi}}\,\eu^{-t^2/2}
\]
(وهي \omterm{def:g10:functions:graph}{منحنى} الجرس). وبوجه أعم،
$X \sim \mathcal N(\mu, \sigma^2)$ إذا كان $\dfrac{X - \mu}{\sigma} \sim
\mathcal N(0,1)$؛ وعندئذٍ $\E(X) = \mu$ و $\V(X) = \sigma^2$.
\end{definition}

\begin{omfigure}
\begin{tikzpicture}
\begin{axis}[
  width=11cm, height=4.6cm,
  axis lines=middle, xmin=-4, xmax=4, ymin=0, ymax=0.45,
  xtick={-3,-2,-1,0,1,2,3},
  xticklabels={$\mu{-}3\sigma$,$\mu{-}2\sigma$,$\mu{-}\sigma$,$\mu$,
               $\mu{+}\sigma$,$\mu{+}2\sigma$,$\mu{+}3\sigma$},
  ytick=\empty, samples=120, domain=-4:4]
  \addplot[name path=gauss, omDef, thick] {exp(-x^2/2)/sqrt(2*pi)};
  \path[name path=xaxis] (-4,0) -- (4,0);
  \addplot[omDef!20] fill between[of=gauss and xaxis,
           soft clip={domain=-2:2}];
  \node[omDef] at (axis cs:0,0.17) {$\approx 95.4\%$};
\end{axis}
\end{tikzpicture}

{\small \omterm{def:g12:contdist:density}{الكثافة} الطبيعية: نحو $68.3\%$ من الكتلة يقع ضمن
$\sigma$ من \omterm{def:g10:stats:mean}{المتوسط}، و $95.4\%$ ضمن $2\sigma$، و $99.7\%$ ضمن
$3\sigma$.}
\end{omfigure}

\begin{remark}
العامل $\frac{1}{\sqrt{2\pi}}$ يجعل المساحة الكلية $1$ --- وهو حساب
شهير (\emph{\omterm{def:g12:integ:area}{تكامل} غاوس}) يُجرى في الجامعة. ولا توجد
صيغة ابتدائية للمقدار $\int \varphi$: فالاحتمالات الطبيعية تُقرأ من
جداول أو آلات حاسبة.
\end{remark}

\begin{theorem}[دي موافر--لابلاس، مقبولة]\label{thm:g12:contdist:tcl}
\index{مبرهنة دي موافر--لابلاس}
لتكن $X_n \sim \mathcal B(n, p)$ ولتكن
$Z_n = \dfrac{X_n - np}{\sqrt{np(1-p)}}$ (وهي الثنائية المعيَّرة). عندئذٍ
يكون من أجل كل $a \leq b$،
\[
\P(a \leq Z_n \leq b) \xrightarrow[n\to+\infty]{}
\int_a^b \varphi(t)\,\dd t .
\]
\emph{وهذه النتيجة مقبولة في هذا المستوى.} وهي تفسّر الظهور
الكوني \omterm{def:g10:functions:graph}{لمنحنى} الجرس: \omterm{def:g12:randvar:binomial}{فالتوزيع الثنائي} ذو $n$ الكبير طبيعي
\omterm{def:g10:numbers:approx}{تقريبًا} --- وكذلك (بمبرهنة النهاية المركزية، في الجامعة) أي مجموع لآثار
صغيرة مستقلة كثيرة.
\end{theorem}

\begin{method}[مجال التذبذب عند $95\%$]\label{met:g12:contdist:fluctuation}
\index{مجال ثقة}
بما أن متغيرًا طبيعيًا يقع ضمن $1.96$ انحرافًا معياريًا من متوسطه
\omterm{def:g10:proba:distribution}{باحتمال} $0.95$، فإنه من أجل $n$ كبير يحقق \omterm{def:g10:stats:series}{تواتر}
$F_n = \frac{X_n}{n}$ \omterm{def:g10:proba:sample}{لعينة} $\mathcal B(n,p)$
\[
\P\left(p - 1.96\sqrt{\tfrac{p(1-p)}{n}} \leq F_n \leq
p + 1.96\sqrt{\tfrac{p(1-p)}{n}}\right) \approx 0.95 .
\]
\begin{itemize}
  \item \emph{مجال التذبذب} (للاختبار): إذا وضع $p$ مدَّعى
        \omterm{def:g10:stats:series}{التواتر} الملاحَظ خارج هذا المجال، فارفض الادعاء عند
        المستوى $5\%$.
  \item \emph{\omterm{met:g12:contdist:fluctuation}{مجال الثقة}} (للتقدير): المجال المبسَّط
        $\left[F_n - \frac{1}{\sqrt n},\ F_n + \frac{1}{\sqrt n}\right]$
        يحوي $p$ \omterm{def:g10:proba:distribution}{باحتمال} لا يقل عن $0.95$
        (باستعمال $1.96\sqrt{p(1-p)} \leq 1$، انظر \cref{exo:g12:sums:8}).
\end{itemize}
\end{method}

\begin{example}\label{ex:g12:contdist:poll}
استطلاع شمل $n = 1000$ ناخب يعطي مرشحًا $52\%$. \omterm{met:g12:contdist:fluctuation}{ومجال الثقة}
$52\% \pm \frac{1}{\sqrt{1000}} \approx 52\% \pm 3.2\%$ ما زال
يحوي $50\%$: فالاستطلاع وحده لا يثبت أن المرشح
متقدم.
\end{example}

\section{تمارين}

\begin{exercise}[$\star$]\label{exo:g12:contdist:1}
تمر حافلة كل $15$ دقيقة؛ ويصل مسافر في لحظة عشوائية
منتظمة. وليكن $X \sim \mathcal U\intcc{0}{15}$ زمن الانتظار. احسب
$\P(X \leq 5)$ و $\P(4 \leq X \leq 10)$ و $\E(X)$ و $\sigma(X)$.
\end{exercise}

\begin{exercise}[$\star$]\label{exo:g12:contdist:2}
تحقق من أن $f(t) = \frac{3}{8}t^2$ \omterm{def:g12:contdist:density}{كثافة} على $\intcc{0}{2}$، ثم
احسب $\E(X)$ و $\P(X \geq 1)$ من أجل متغير $X$ له هذه \omterm{def:g12:contdist:density}{الكثافة}.
\end{exercise}

\begin{exercise}[$\star$]\label{exo:g12:contdist:3}
عمر مكوّن إلكتروني (بالسنوات) يتبع
$\mathcal E(0.2)$.
\begin{enumerate}
  \item احسب العمر المأمول و $\P(X > 5)$.
  \item عمل المكوّن $3$ سنوات أصلًا. فما \omterm{def:g10:proba:distribution}{احتمال}
        أن يعمل $5$ سنوات أخرى على الأقل؟
\end{enumerate}
\end{exercise}

\begin{exercise}[$\star\star$]\label{exo:g12:contdist:4}
نصف عمر ذرة مشعّة عمرها يتبع
$\mathcal E(\lambda)$ هو \omterm{def:g11:stat:median}{الوسيط} $m$: $\P(X > m) = \frac12$. عبّر عن
$m$ بدلالة $\lambda$ وقارنه \omterm{def:g12:randvar:exp}{بالأمل الرياضي}. فأيهما
أكبر، ولماذا يكون ذلك معقولًا من أجل \omterm{def:g12:randvar:rv}{توزيع} غير متناظر؟
\end{exercise}

\begin{exercise}[$\star\star$]\label{exo:g12:contdist:5}
أطوال في مجتمع تتبع $\mathcal N(175,\ 7^2)$ (بالسنتيمتر). وباستعمال
قاعدة $68$--$95$--$99.7$، قدّر نسبة السكان الذين يتراوح
طولهم بين $168$ و $182$ cm، وفوق $189$ cm، ودون $154$ cm.
\end{exercise}

\begin{exercise}[$\star\star$]\label{exo:g12:contdist:6}
آلة تملأ أكياسًا مكتوبًا عليها $500$ g؛ والكتلة المعبَّأة تتبع
$\mathcal N(\mu,\ 4^2)$. وتقتضي الأنظمة ألا يزيد عدد الأكياس التي
تزن دون $500$ g على $2.5\%$. وباستعمال قاعدة $95\%$، ما أدنى ضبط
للمقدار $\mu$ يمتثل لذلك؟
\end{exercise}

\begin{exercise}[$\star\star$]\label{exo:g12:contdist:7}
حجر نرد يُدَّعى أنه غير مغشوش يُرمى $1200$ مرة فيُظهر ستة $260$
مرة (كما في \cref{exo:g12:sums:7}).
\begin{enumerate}
  \item احسب مجال التذبذب عند $95\%$ \omterm{def:g10:stats:series}{لتواتر} الأوجه الستة
        من أجل حجر غير مغشوش على $1200$ رمية.
  \item فهل \omterm{def:g10:stats:series}{التواتر} الملاحَظ داخله؟ وقارن قوة هذا
        الاستنتاج بتحليل بيينيمي--تشيبيشيف.
\end{enumerate}
\end{exercise}

\begin{exercise}[$\star\star\star$]\label{exo:g12:contdist:8}
قبل انتخاب، سيقدّر استطلاع يشمل $n$ شخصًا نتيجة $p$ لمرشح
بواسطة \omterm{def:g10:stats:series}{التواتر} الملاحَظ $F_n$.
\begin{enumerate}
  \item \omterm{met:g12:contdist:fluctuation}{بمجال الثقة} في
        \cref{met:g12:contdist:fluctuation}، أي حجم \omterm{def:g10:proba:sample}{عينة} يضمن
        هامشًا قدره $\pm 2$ نقطة؟
  \item يفصل بين المرشحين في الواقع $1$ نقطة. فسّر لماذا
        لا يستطيع أي استطلاع واقعي أن يحسم الفائز حسمًا موثوقًا، مهما
        أُحسن إجراؤه.
\end{enumerate}
\end{exercise}

\begin{exercise}[$\star\star\star$]\label{exo:g12:contdist:9}
ليكن للمتغير $X$ \omterm{def:g12:contdist:density}{الكثافة} $f$ على $\intcc ab$ وليكن
$Y = \alpha X + \beta$ حيث $\alpha > 0$.
\begin{enumerate}
  \item بيّن أن $\P(c \leq Y \leq d)
        = \P\left(\frac{c - \beta}{\alpha} \leq X \leq
        \frac{d-\beta}{\alpha}\right)$، واستنتج أن للمتغير $Y$ \omterm{def:g12:contdist:density}{الكثافة}
        $g(y) = \frac{1}{\alpha} f\left(\frac{y - \beta}{\alpha}\right)$.
  \item استنتج أنه إذا كان $Z \sim \mathcal N(0,1)$، فإن
        $\mu + \sigma Z$ له \omterm{def:g12:contdist:density}{الكثافة}
        $\frac{1}{\sigma\sqrt{2\pi}}\,
        \eu^{-(y - \mu)^2/(2\sigma^2)}$ --- وهي \omterm{def:g12:contdist:density}{الكثافة} الطبيعية العامة.
\end{enumerate}
\end{exercise}

\section{مسألة: في انتظار الحافلة، وقياس الجمهور}

\begin{problem}[{مسألة نهاية الأسبوع --- ثلاث كثافات تدير
العالم: الجهل التام، والانتظار بلا ذاكرة، وجرس الأسباب
الصغيرة الكثيرة؛ زائد المفارقة التي تجعل حافلتك متأخرة دائمًا}]
\label{pb:g12:contdist:1}
لماذا تبدو حافلتك دائمًا أطول انتظارًا مما يعد به جدول
المواعيد؟ ولماذا يكون القسم الذي يجلس فيه التلميذ الوسطي أكبر
من القسم الوسطي؟ ولماذا يكون لأصدقائك أصدقاء أكثر
منك؟ قطعة ماكرة واحدة من الرياضيات --- المعاينة المنحازة
بالطول --- تجيب عن الثلاثة، وهي تعيش في هذا الفصل،
بين التوزيعات المنتظم والأسي والطبيعي
(\cref{def:g12:contdist:uniform} و \cref{def:g12:contdist:expo} و
\cref{def:g12:contdist:normal}). وتشغّل مسألة الكتاب
الأخيرة هذه الكثافات الثلاث وتنتهي حيث تشير السلسلة
كلها: إلى \omterm{def:g10:functions:graph}{منحنى} الجرس.

\medskip
\noindent\textbf{الجزء الأول --- ثلاث شخصيات.}
\begin{enumerate}
  \item تصل حافلة في لحظة عشوائية منتظمة خلال
        $30$ دقيقة القادمة. احسب $\P(\text{الانتظار} \leq 10)$، و
        الانتظار المأمول، و $\sigma$
        (\cref{prop:g12:contdist:uniformmoments}).
  \item تحقق من أن $f(x) = 2x$ على $\intcc{0}{1}$ \omterm{def:g12:contdist:density}{كثافة}،
        واحسب $\E(X)$ و $\P\left(X > \frac12\right)$.
  \item تصل مكالمات هاتفية إلى خط مساعدة تدفقًا بلا ذاكرة:
        وانتظار المكالمة التالية أسي \omterm{def:g10:stats:mean}{بمتوسط}
        $10$ دقائق. احسب $\P(T > 15)$ والانتظار \omterm{def:g11:stat:median}{الوسيط}
        --- ولماذا يكون \omterm{def:g11:stat:median}{الوسيط} \emph{أقل} من \omterm{def:g10:stats:mean}{المتوسط}؟
  \item وقد انتظرت $20$ دقيقة أصلًا. فكم يساوي
        $\P(T > 35 \mid T > 20)$
        (\cref{thm:g12:contdist:memoryless})؟ فسّر في جملة
        واحدة.
  \item طابق كل انتظار بنموذجه الصحيح --- منتظم،
        أو أسي، أو لا هذا ولا ذاك: (أ) نقرة عداد
        غايغر؛ (ب) مترو يمر كل $8$ دقائق، ووصولك
        غير متزامن؛ (ج) المكالمة التالية على خط
        المساعدة؛ (د) عطب مصباح كهربائي \emph{يبلى}.
        برّر (د) بمبرهنة انعدام الذاكرة.
\end{enumerate}

\medskip
\noindent\textbf{الجزء الثاني --- الجرس في العمل.}
(استعمل معالم \omterm{def:g12:contdist:normal}{التوزيع الطبيعي}: $68\,\%$ ضمن $\sigma$، و $95\,\%$
ضمن $2\sigma$، و $99.7\,\%$ ضمن $3\sigma$.)
\begin{enumerate}[resume]
  \item أطوال البالغين تتبع \omterm{def:g10:numbers:approx}{تقريبًا}
        $\mathcal N(172,\ 8)$ (بالسنتيمتر). فما النسب الواقعة في
        $\intcc{164}{180}$ و $\intcc{156}{188}$ و
        $\intcc{148}{196}$؟
  \item قدّر نسبة من يزيد طولهم على $190$ cm (\omterm{pb:g11:stat:1}{بالدرجة المعيارية}
        أولًا؛ ثم بجدول أو آلة حاسبة).
  \item درجات الذكاء معايَرة على $\mathcal N(100, 15)$.
        فأي نسبة تسجّل فوق $130$ --- وهي نحو شخص
        واحد من كم؟
  \item آلة تقطع براغي أطوالها
        $\mathcal N(50,\ 0.1)$ mm والمواصفة هي
        $\intcc{49.8}{50.2}$. فما الحصة المرفوضة؟ وتحتفي الصناعة
        بعمليات ``ستة سيغما''، التي تجلس
        مواصفاتها عند $\pm 6\sigma$: فماذا يشتري ذلك،
        ولماذا تسعى المصانع وراءه؟
  \item مبرهنة دي موافر--لابلاس (\cref{thm:g12:contdist:tcl}): من أجل
        $100$ رمية لقطعة نقد غير مغشوشة، قرّب
        $\P(45 \leq \text{صورة} \leq 55)$ (مع تصحيح \omterm{def:g12:limcont:continuity}{الاتصال}:
        $\pm 0.5$؛ و $\sigma = 5$). فأي آلة خشبية في
        \cref{pb:g11:binom:1} تملّسها هذه المبرهنة إلى
        \omterm{def:g10:functions:graph}{منحنى}؟
  \item انطلاقًا من \cref{exo:g12:contdist:8}: يحتاج هامش $\pm 2$ نقطة
        إلى نحو $n = 2\,400$. احسب \omterm{def:g10:proba:sample}{العينة} اللازمة من أجل
        $\pm 0.5$ نقطة --- واختم، في جملة واحدة، لماذا
        لا يستطيع أي استطلاع أن يحسم بأمانة سباقًا يفصل بينه نقطة واحدة.
\end{enumerate}

\medskip
\noindent\textbf{الجزء الثالث --- مفارقة الحافلة.}
\begin{enumerate}[resume]
  \item تصل الحافلات تدفقًا بلا ذاكرة \omterm{def:g10:stats:mean}{متوسط} فجوته $10$
        دقائق. وتبلغ أنت الموقف في لحظة كيفية.
        ويقول الحدس إن الانتظار المأمول $5$ دقائق
        (أي نصف فجوة). فماذا تقول مبرهنة انعدام الذاكرة
        بدل ذلك؟
  \item شخّص الحدس بجدول مواعيد لعبة: فجوات
        تتناوب $5$ و $15$ دقيقة (\omterm{def:g10:stats:mean}{ومتوسط} الفجوة $10$).
        احسب \omterm{def:g10:proba:distribution}{احتمال} أن يقع وصول عشوائي منتظم
        في فجوة طويلة، ثم الانتظار المأمول الحقيقي ---
        وسمِّ الجاني: فالوصولات تعاين الفجوات \omterm{def:g10:proba:distribution}{باحتمال}
        يتناسب مع \emph{طولها}.
  \item وفي التدفق عديم الذاكرة، للفجوة التي تقع فيها طول
        مأمول قدره $20$ دقيقة --- أي ضعف الفجوة النموذجية
        (فالزمن إلى الوراء حتى الحافلة \omterm{def:g10:functions:function}{السابقة} والزمن إلى الأمام حتى
        التالية \emph{كلاهما} أسي \omterm{def:g10:stats:mean}{بمتوسط} $10$).
        وفّق بين هذا والسؤال 12 وصُغ
        \emph{مفارقة المعاينة} في جملة واحدة.
  \item اللسعة نفسها على اليابسة: جامعة فيها تسعة أقسام من
        $10$ طلاب وقسم من $110$. احسب \omterm{def:g10:stats:mean}{متوسط}
        حجم القسم، ثم حجم القسم الذي يعيشه \emph{الطالب
        الوسطي}. فأي عدد سيقتبسه الكتيّب،
        وأيهما حقيقة الطلاب المعيشة؟
  \item ومفارقة الصداقة --- ``لأصدقائك، وسطيًا، أصدقاء أكثر
        منك'' --- هي الرياضيات نفسها. قل في جملة واحدة ما الذي يلعب دور
        الفجوة الطويلة.
\end{enumerate}

\medskip
\noindent\textbf{الجزء الرابع --- إغلاق الكتاب.}
\begin{enumerate}[resume]
  \item المحاكاة، جسر الممارسين: إذا كان $U$
        منتظمًا على $\intoo{0}{1}$، فبيّن أن
        $T = -10\ln U$ أسي \omterm{def:g10:stats:mean}{بمتوسط} $10$ (بحساب
        $\P(T > t)$). فكل محاكاة انتظار عشوائي في
        الصناعة تعمل بهذه الحيلة ذات السطر الواحد.
  \item وصفة قديمة عند مبرمجي الألعاب تبني جرسًا: اجمع
        اثني عشر متغيرًا منتظمًا مستقلًا على $\intcc{0}{1}$ ثم
        اطرح $6$. أعطِ \omterm{def:g10:stats:mean}{متوسط} المجموع \omterm{def:g12:randvar:exp}{وتباينه}،
        وفسّر --- ولوح غالتون في يدك --- لماذا تكون
        النتيجة طبيعية \omterm{def:g10:numbers:approx}{تقريبًا}.
  \item الأسواق تنهار أشد مما يسمح به الجرس: فقد وُصف انهيار
        سنة 1987 بأنه ``\omterm{def:g11:prob:model}{حادثة} $25\sigma$''، وهذه
        احتمالها في ظل الطبيعية نحو $10^{-137}$.
        فما الاستنتاج الصحيح --- عن العالم، أم
        عن النموذج؟ (جملة واحدة، وهي عقيدة الإحصائي.)
  \item الخاتمة، ووداع الكتاب: الشخصيات الثلاث
        في سطر واحد لكل منها --- المنتظم (جهل
        داخل حدود)، والأسي (خطر بلا ذاكرة)، و
        الطبيعي (مجموع أسباب صغيرة كثيرة)؛ ومفارقة المعاينة
        لسعةً للفصل؛ وقوس السلسلة
        كلها --- من عدّ علب العصير في الكتاب السابق إلى
        \omterm{def:g10:functions:graph}{منحنى} الجرس الذي يقيس كل جمهور --- والكتب
        الجامعية تنتظر حيث تتسلم النهايات
        والتكاملات ومبرهنة النهاية المركزية الراية.
\end{enumerate}
\end{problem}
